% GENERATED by docs/manual/build.py - do not edit; change the source it names instead.
% source: docs/pki.md
\hypertarget{doc-pki}{%
\chapter{Certificates / PKI governance}\label{doc-pki}}

\begin{itemize}
\tightlist
\item
  \textbf{Private CA signer}: RSA key with a CA certificate. \texttt{pki.yaml\ -\textgreater{}\ ca\_keys} creates a DFC key with
  a self-signed certificate whose config carries \texttt{basicConstraints=critical,\ CA:TRUE} and
  \texttt{keyUsage=critical,\ keyCertSign,\ digitalSignature,\ cRLSign} (\texttt{create-dfc-key\ -\/-generate-self-signed-certificate\ true\ -\/-conf-file-path\ csr.conf}). An existing CA can be uploaded
  with \texttt{upload-rsa\ -\/-cert} and referenced by name. Classic keys can sign only if Classic is enabled
  as protection-key type in account settings.
\item
  \textbf{Issuer} = the certificate policy: allowed domains/SANs, wildcard/subdomain rules, EKU flags
  (\texttt{server\_flag}, \texttt{client\_flag}, \texttt{code\_signing\_flag}), key usage, subject fields, max TTL.
\item
  \texttt{destination\_path} stores issued certificates as items - required for \texttt{auto\_renew}
  (+ \texttt{scheduled\_renew\_days}) and CRL. \texttt{protect\_certificates} prevents their deletion.
\item
  CRL: \texttt{create\_public\_crl} (served by Akeyless at \texttt{https://vault.akeyless.io/crl/\textless{}account-id\textgreater{}/\textless{}issuer-display-id\textgreater{}})
  or \texttt{create\_private\_crl} (served by the Gateway; needs \texttt{gw\_cluster\_url}, port 8000). OCSP:
  \texttt{create\_public\_ocsp} / \texttt{create\_private\_ocsp}. ACME: \texttt{enable\_acme} (needs \texttt{gw\_cluster\_url}).
\item
  \textbf{Public CA} (Let\textquotesingle s Encrypt, ZeroSSL, DigiCert, Sectigo, GlobalSign, GlobalSign Atlas, Google
  Trust...): create the CA target (\texttt{target\ create\ \textless{}ca\textgreater{}} / \texttt{akeyless\_target\_\textless{}ca\textgreater{}}), then an issuer with
  \texttt{ca\_target} + \texttt{gw\_cluster\_url}. Sectigo approval workflows are not supported by the issuer.
\item
  Expiry notifications: \texttt{expiration\_event\_in\_days} on the issuer + an event forwarder with
  \texttt{certificate-pending-expiration} / \texttt{certificate-expired}.
\item
  \textbf{RBAC}: issuers are items. Requesters need a rule on the issuer and on its destination path.
  The exact capability \texttt{get-pki-certificate} requires is not documented - \textbf{VERIFY} with
  \texttt{describe-permissions} (start with \texttt{read}; the reference model grants \texttt{read,list}).
\end{itemize}
